\section{NVIDIA is a simulation company}

上一章讲产业链分工，本章聚焦 NVIDIA。谢晨回忆黄仁勋在内部说过：NVIDIA is a simulation company。这句话的意义不只是营销，而是把 NVIDIA 从“卖 GPU”重新理解为“用计算创造可训练世界”的公司。游戏是服务人类体验的仿真，Omniverse 是服务工业和数字孪生的仿真，Isaac Sim 是服务机器人和物理 AI 的仿真。

NVIDIA 的 three computer problem 也连接了这一点：第一个计算机是数据中心，第二个是端侧物理 AI 设备，第三个是仿真计算机。数据中心训练大模型，端侧机器人和车运行模型，仿真计算机生成物理世界、数据和评估环境。机器人数据荒越严重，第三台计算机越像战略入口。

\begin{figure}[H]
\centering
\includegraphics[width=0.96\textwidth]{nvidia-simulation-company.png}
\caption{NVIDIA is a simulation company：算力、物理引擎、仿真环境和机器人训练连成闭环。自制概念图，依据 01:15:33--01:21:25 对谈内容整理。}
\end{figure}

\begin{knowledgebox}{读图：NVIDIA 的仿真战略连接三类计算}
图中左侧是训练和推理算力，中间是物理引擎、渲染、资产和 API，右侧是机器人与自动驾驶部署。NVIDIA 的强处不是单点工具，而是把 GPU、Omniverse、Isaac Sim、PhysX、Warp/Newton 和机器人生态连成基础设施。仿真越成为数据生产方式，算力公司越接近数据公司。
\end{knowledgebox}

\subsection{Isaac Sim、MuJoCo 与下一代物理引擎}

本节把访谈里的仿真器生态稍微结构化。NVIDIA 的 Isaac Sim 建在 PhysX 和 Omniverse 渲染之上，优势是 GPU、渲染和工业生态；Google/DeepMind 相关的 MuJoCo 长期是机器人研究常用物理引擎，优势是物理和 API，但渲染不是强项。随着端到端 RL 需要大量 GPU 并行，CPU 版 MuJoCo 难以满足吞吐，MJX 和 MuJoCo Warp/Newton 代表了向 GPU 加速、开源生态和统一物理底座演进的方向。

这里最值得学习的不是某个引擎谁赢，而是仿真器会成为生态竞争的一层。底层物理引擎越开源、越高效、越可维护，上层资产、API、数据生成和 RL 平台越能快速发展。光轮这样的公司如果能贡献资产、物理参数和工具链，就会成为这个生态中的重要部分，而不一定要独占底层引擎。

\begin{knowledgebox}{术语消化：仿真栈四层}
\begin{tabularx}{\textwidth}{p{0.18\textwidth}p{0.36\textwidth}X}
\toprule
层级 & 代表内容 & 作用 \\
\midrule
Physics Solver & PhysX、MuJoCo、Newton 等 & 计算力学、碰撞、约束和运动。 \\
Rendering & Omniverse 等渲染管线 & 生成视觉观测和传感器输入。 \\
Sim-ready Assets & 可交互资产、场景、机器人本体 & 提供训练世界的原料。 \\
API/Framework & 数据导出、metadata、RL 环境接口 & 让仿真能被模型训练和评价调用。 \\
\bottomrule
\end{tabularx}
\end{knowledgebox}

\subsection{合成数据与大模型}

本节从机器人扩展到大模型。访谈中谢晨判断，自然语言大模型也越来越依赖合成数据，因为互联网数据趋于穷尽，模型本身可以生成新任务、新解法和新评价。某种意义上，GPT 也是合成数据生成器，它不断产生自然语言并可被用来训练、蒸馏或评估其他模型。

对具身大模型而言，合成数据则更复杂：它不仅需要文本，还需要 3D、机器人视角、多传感器、动作轨迹和物理交互。自然语言合成数据可以依赖模型生成和验证，具身合成数据必须依赖仿真世界和真实物理校准。这也是为什么 EP109 把合成数据讲得比一般大模型讨论更重。

\begin{warningbox}{语言合成数据和具身合成数据不是同一难度}
语言数据可以主要在符号空间中生成和检查；具身数据必须经过物理世界约束。一个机器人轨迹是否有效，不能只看文本逻辑，还要看接触、力学、安全、执行器和真实部署。
\end{warningbox}

\subsection{本章小结}

NVIDIA 的仿真战略说明，未来 AI 基础设施不只有训练集群和端侧芯片，还包括生成训练世界的仿真计算机。对机器人而言，仿真既是数据生产工具，也是模型评价环境和产业生态入口。

